\documentclass[12pt]{article}

\usepackage{graphicx}
\usepackage{amsmath}
\usepackage[colorlinks]{hyperref}
\usepackage{endfloat}

\title{Problema 6}
\author{Joaquín Rapela}

\begin{document}

\maketitle

\section*{Parte a}

La expresión del campo eléctrico sobre el eje $x$ es:

\begin{align}
    E(x, 0, 0)&=\frac{q}{4\pi\epsilon_0}\left(\frac{x+d/2}{|x+d/2|^3}-\frac{x-d/2}{|x-d/2|^3}\right)\hat{i}
\end{align}

La figura~\ref{fig:6a1_3} grafica $E(x, 0, 0)$ en función de $x$.  El sentido
del $E(x, 0, 0)$ varía a lo largo del eje $x$. La figura~\ref{fig:6a5} grafica
$||E(x, 0, 0)||_2$ en función de $x$.


\begin{figure}
    \centering
    \href{https://www.gatsby.ucl.ac.uk/~rapela/em/delfi/mySolutions/guia1/figures/6a1_3.html}{\includegraphics[width=6in]{../figures/6a1_3.png}}

    \caption{E(x, 0, 0) para $q=0.1\mu C$, y $d=8 m$. Haga click en la imagen
    para ver su version interactiva. El código usado para construir esta figura
    se encuentra
    \href{https://github.com/joacorapela/em/blob/master/delfi/mySolutions/guia1/code/scripts/6a1_3.py}{aquí}.}

    \label{fig:6a1_3}
\end{figure}

\begin{figure}
    \centering
    \href{https://www.gatsby.ucl.ac.uk/~rapela/em/delfi/mySolutions/guia1/figures/6a5.html}{\includegraphics[width=6in]{../figures/6a5.png}}

    \caption{$||E(x, 0, 0)||_2$ para $q=0.1\mu C$, y $d=8 m$. Haga click en la
    imagen para ver su version interactiva. El código usado para construir esta
    figura se encuentra
    \href{https://github.com/joacorapela/em/blob/master/delfi/mySolutions/guia1/code/scripts/6a5.py}{aquí}.}

    \label{fig:6a5}
\end{figure}

\section*{Parte b}

La expresión del campo eléctrico sobre el eje $y$ es:

\begin{align}
    E(0, y, 0)&=\frac{qd}{(4\pi\epsilon_0)((d/2)^2+y^2)^{3/2}}\hat{i}
\end{align}

La figura~\ref{fig:6b} grafica $E(0, y, 0)$ en función de $y$.  El sentido
del $E(0, y, 0)$ no varía a lo largo del eje $y$.


\begin{figure}
    \centering
    \href{https://www.gatsby.ucl.ac.uk/~rapela/em/delfi/mySolutions/guia1/figures/6b.html}{\includegraphics[width=6in]{../figures/6b.png}}

    \caption{E(0, y, 0) para $q=0.1\mu C$, y $d=8 m$. Haga click en la imagen
    para ver su version interactiva. El código usado para construir esta figura
    se encuentra
    \href{https://github.com/joacorapela/em/blob/master/delfi/mySolutions/guia1/code/scripts/6b.py}{aquí}.}
    \label{fig:6b}
\end{figure}

\section*{Parte c}

La expresión del campo eléctrico sobre el eje $y$ es:

\begin{align*}
    E(\mathbf{p})&=\frac{q}{4\pi\epsilon_0}
    \left(
    \frac{\mathbf{p}-\begin{bmatrix}
                       -d/2\\
                       0\\
                       0
                     \end{bmatrix}}
         {\left\Vert\mathbf{p}-\begin{bmatrix}
                                 -d/2\\
                                 0\\
                                 0
                               \end{bmatrix}\right\Vert^3}-
    \frac{\mathbf{p}-\begin{bmatrix}
                       d/2\\
                       0\\
                       0
                     \end{bmatrix}}
         {\left\Vert\mathbf{p}-\begin{bmatrix}
                                 d/2\\
                                 0\\
                                 0
                               \end{bmatrix}\right\Vert^3}
    \right)
\end{align*}

La figura~\ref{fig:6c} grafica $E(\mathbf{p})$ en el espacio.


\begin{figure}
    \centering
    \href{https://www.gatsby.ucl.ac.uk/~rapela/em/delfi/mySolutions/guia1/figures/6c.html}{\includegraphics[width=6in]{../figures/6c.png}}

    \caption{$E(\mathbf{p})$ para $q=0.1\mu C$, y $d=8 m$. Haga click en la imagen
    para ver su version interactiva. El código usado para construir esta figura
    se encuentra
    \href{https://github.com/joacorapela/em/blob/master/delfi/mySolutions/guia1/code/scripts/6c.py}{aquí}.}
    \label{fig:6c}
\end{figure}

\end{document}
